\bookchapter{}{The mathematical toolbox}{Every tool the lessons use beyond a first course in calculus, differential equations and linear algebra — each derived once, in one place.}
\label{app:math}

\lead{The lessons introduce their mathematics where they need it, and only as much as they need. That is the right order for learning control and the wrong one for looking something up. The appendices that follow collect the same tools by subject. Each is a short course: definitions, theorems with a proof or a source, and worked examples with the numbers of the drone. None of it is required before Lesson~1. All of it will have been used by the end of the book.}

\section*{What is in it}

\begin{center}\small
\renewcommand{\arraystretch}{1.3}
\begin{tabularx}{\linewidth}{@{}>{\sffamily\bfseries\color{accent}}l>{\raggedright\arraybackslash}X>{\raggedright\arraybackslash}p{30mm}@{}}
\toprule
 & \normalfont\sffamily\color{ink} subject & \sffamily first needed in\\
\midrule
A & \hyperref[app:ode]{Linear differential equations and state space}: characteristic equation, the matrix exponential, controllability & \lessonref{les:spring}\\
B & \hyperref[app:routh]{Polynomials and stability}: the Routh–Hurwitz criterion, the root locus & \lessonref{les:integral}\\
C & \hyperref[app:laplace]{Complex numbers and transforms}: frequency response, Fourier series, the Laplace transform, transfer functions, the Nyquist criterion & \lessonref{les:noise}\\
D & \hyperref[app:sampling]{Sampling and the z-transform}: aliasing, difference equations, exact discretisation & \lessonref{les:slow}\\
E & \hyperref[app:prob]{Probability and random signals}: variance and covariance, white noise, spectra, least squares & \lessonref{les:noise}\\
F & \hyperref[app:num]{Numerical methods}: roots, integrals, the integration of differential equations & \lessonref{les:meet}\\
G & \hyperref[app:stab]{Stability of nonlinear systems}: linearisation, Lyapunov functions, the invariance principle, time scales & \lessonref{les:spring}\\
H & \hyperref[app:quat]{Rotations and quaternions}: rotation matrices, axis and angle, quaternion kinematics, the attitude error & \lessonref{les:tilt}\\
\bottomrule
\end{tabularx}
\end{center}

Two more appendices belong to Part~II and are printed with it: \textsf{\textbf{I}}, linear algebra beyond eigenvalues (singular values, positive definite matrices, the Riccati equation), and \textsf{\textbf{J}}, optimisation (convexity, the Karush–Kuhn–Tucker conditions, quadratic and cone programs). The later volumes add \textsf{\textbf{K}}, norms of signals and systems, and \textsf{\textbf{L}}, the calculus of variations and dynamic programming. The letters stay the same throughout the course.

\section*{How to use it}

\begin{itemize}
  \item \textbf{As a reference.} Each appendix opens with a margin card, \emph{Where it is used}, that names the lessons relying on it, and the lessons point back here. Equations and theorems are numbered with the letter of the appendix: \enquote{Theorem~A.2}, \enquote{Equation~(C.7)}.
  \item \textbf{As a refresher.} The appendices can be read in order. A and C together are the classical \enquote{signals and systems} course in forty pages; B and G are its stability theory; D, E and F are what a digital, noisy, computed loop adds; H is the geometry of a body that turns.
  \item \textbf{With a pencil.} The worked examples use the drone of the lessons — $m = \SI{1}{kg}$, $K_p = 10$, $K_d = 7$, $K_i = 0.8$, a motor lag of \SI{30}{ms}, a controller at \SI{250}{Hz} — so that every number can be checked against a chapter or against the simulator.
\end{itemize}

The statements follow the conventions of the \hyperref[sec:notation]{Prelude}: a \emph{Definition} fixes a term, a \emph{Theorem} is general and comes with a proof or with the source of one, a \emph{Result} is a formula for our vehicle. The toolbox has no exercises; those are in the lessons.

\toolchapter{A}{Linear differential equations and state space}{One guess, $e^{st}$, solves them all. The rest is bookkeeping — and the bookkeeping has a name: the state.}
\label{app:ode}

\lead{\usedcard{\setuprow{Lesson 2}{the characteristic equation, poles}\setuprow{Lessons 3–5}{third-order loops, fitting initial conditions}\setuprow{Lessons 1, 9}{the state, the matrix exponential, the sampled loop}\setuprow{Lessons 4, 12}{controllability}\setuprow{Part II}{state feedback, observers: everything}}%
Every closed loop in Part~I is a linear differential equation with constant coefficients, or is treated as one. This appendix solves such equations once and for all: first the scalar ones, of first, second and any order; then the same equations written as a system of first-order ones, $\dot{\vx} = A\vx + B\symbf u$, whose solution is a single formula with a matrix in the exponent.}

\section{First order: the exponential and its time constant}

The simplest dynamical system there is has one variable whose rate of change is proportional to its distance from a target:
\begin{equation}\label{eq:A-first}
  \tau\,\dot x + x = u , \qquad \tau > 0 .
\end{equation}
The motors of the drone obey it ($x$ is the thrust, $u$ the command), and so do the D filter of \lessonref{les:noise}, the creeping integral of \lessonref{les:droop} and each loop of the cascade in \lessonref{les:cascade}.

\paragraph{No input.} For $u = 0$ the equation reads $\dot x = -x/\tau$. Separating the variables, $\dd x/x = -\dd t/\tau$, and integrating from $0$ to $t$ gives $\ln x(t) - \ln x_0 = -t/\tau$, that is,
\begin{equation}
  x(t) = x_0\,e^{-t/\tau} .
\end{equation}
The number $\tau$ is the \term{time constant}. After one time constant $x$ has fallen to $e^{-1} = 37\,\%$ of its initial value, after three to 5\,\%, after five to below 1\,\%.\tip{The tangent to $e^{-t/\tau}$ at any moment reaches zero exactly $\tau$ later. On a chart, draw the tangent at the start of a response: where it meets the final value, read off $\tau$.}

\paragraph{A constant input.} For $u = u_0$ the constant $x = u_0$ is a solution. Subtract it: $z = x - u_0$ obeys $\tau\dot z + z = 0$, which we have just solved. Hence
\begin{equation}\label{eq:A-step}
  x(t) = u_0 + (x_0 - u_0)\,e^{-t/\tau} .
\end{equation}
The response starts at $x_0$, ends at $u_0$ and covers 63\,\% of the distance in the first time constant.

\begin{example}[A motor spins up]
\label{ex:A-motor}
A motor of the drone delivers \SI{2.45}{N} in hover. The command jumps to \SI{3.45}{N}. With $\tau_m = \SI{30}{ms}$, what is the thrust after one controller period (\SI{4}{ms}), after \SI{30}{ms} and after \SI{90}{ms}?

\textbf{Step 1 — identify.} $x_0 = 2.45$, $u_0 = 3.45$, so $x_0 - u_0 = -1$ and $x(t) = 3.45 - e^{-t/0.03}$.

\textbf{Step 2 — evaluate.} $e^{-4/30} = 0.875$: after \SI{4}{ms} the thrust is \SI{2.575}{N}, 12.5\,\% of the way. $e^{-1} = 0.368$: after \SI{30}{ms}, \SI{3.08}{N}. $e^{-3} = 0.050$: after \SI{90}{ms}, \SI{3.40}{N}.

\textbf{Step 3 — read it.} A controller that runs every \SI{4}{ms} sees only an eighth of each of its commands carried out before it issues the next. That is why the kick of \lessonref{les:kick} hardly reaches the drone.
\end{example}

\paragraph{Any input.} Multiply \cref{eq:A-first} by $e^{t/\tau}/\tau$. The left-hand side becomes the derivative of a product, $\frac{\dd}{\dd t}\big(e^{t/\tau}x\big) = e^{t/\tau}u/\tau$, and integrating gives
\begin{equation}\label{eq:A-conv1}
  x(t) = e^{-t/\tau}\,x_0 + \frac{1}{\tau}\int_0^t e^{-(t - s)/\tau}\,u(s)\,\dd s .
\end{equation}
The first term is what the initial value leaves behind. The second adds up the input of every past moment $s$, each weighted by how much of it survives until $t$. Everything in this appendix is a generalisation of this one line.

\section{Second order: the characteristic equation}

A mass with a spring and a damper — the P–D loop of \lessonref{les:damper} — obeys
\begin{equation}\label{eq:A-second}
  m\,\ddot x + c\,\dot x + k\,x = 0 .
\end{equation}
Try $x = e^{st}$. Each derivative multiplies by $s$, so the equation becomes $(ms^2 + cs + k)\,e^{st} = 0$, and since the exponential never vanishes,
\begin{equation}
  m s^2 + c s + k = 0 .
\end{equation}
This \term{characteristic equation} has two roots, $s_{1,2} = \big({-c} \pm \sqrt{c^2 - 4mk}\,\big)/2m$. Three cases arise, and \lessonref{les:damper} gave them their names.

\begin{result}[The general solution of \cref{eq:A-second}]
\begin{center}\small
\renewcommand{\arraystretch}{1.35}
\begin{tabular}{@{}lll@{}}
\toprule
roots & general solution & name\\
\midrule
real, $s_1 \ne s_2$ & $C_1 e^{s_1 t} + C_2 e^{s_2 t}$ & overdamped\\
real, $s_1 = s_2 = s$ & $(C_1 + C_2\,t)\,e^{st}$ & critically damped\\
complex, $\sigma \pm j\omega$ & $e^{\sigma t}\,(C_1\cos\omega t + C_2\sin\omega t)$ & underdamped\\
\bottomrule
\end{tabular}
\end{center}
The two constants are fixed by the two initial conditions $x(0)$ and $\dot x(0)$.
\end{result}

Where do the second and third lines come from? For complex roots, $e^{(\sigma \pm j\omega)t} = e^{\sigma t}(\cos\omega t \pm j\sin\omega t)$ by Euler's formula; the sum and the difference of the two solutions are real, and they are the cosine and the sine of the table. For a double root only one exponential is available, and a second solution is needed. Try $x = t\,e^{st}$: then $\dot x = (1 + st)e^{st}$ and $\ddot x = (2s + s^2t)e^{st}$, so
\begin{equation*}
  m\ddot x + c\dot x + kx = \big[(ms^2 + cs + k)\,t + (2ms + c)\big]e^{st} .
\end{equation*}
The first bracket vanishes because $s$ is a root, the second because it is a \emph{double} root: $2ms + c$ is the derivative of the characteristic polynomial, and a polynomial and its derivative vanish together exactly at a repeated root.

\begin{example}[The heavy drone, step by step]
\label{ex:A-heavy}
The \SI{1.4}{kg} drone of \lessonref{les:integral} flies with $K_p = 10$, $K_d = 7$ and exact feedforward. It is \SI{1}{m} below its setpoint and at rest. Find its motion.

\textbf{Step 1 — the equation.} $1.4\,\ddot x + 7\,\dot x + 10\,x = 0$ with $x(0) = -1$, $\dot x(0) = 0$.

\textbf{Step 2 — the roots.} $c^2 - 4mk = 49 - 56 = -7$: complex. $s = (-7 \pm j\sqrt7)/2.8 = -2.5 \pm 0.945j$. The extra mass has made the default tune, which is overdamped for \SI{1}{kg}, slightly underdamped.

\textbf{Step 3 — the general solution.} $x(t) = e^{-2.5t}(C_1\cos 0.945t + C_2\sin 0.945t)$.

\textbf{Step 4 — the first condition.} $x(0) = C_1 = -1$.

\textbf{Step 5 — the second condition.} Differentiate with the product rule and set $t = 0$: $\dot x(0) = -2.5\,C_1 + 0.945\,C_2 = 0$, so $C_2 = 2.5\,C_1/0.945 = -2.646$.

\textbf{Step 6 — the answer.} $x(t) = -e^{-2.5t}(\cos 0.945t + 2.646\sin 0.945t)$. Check: $x(0) = -1$; and in the standard form of \cref{eq:damper-std}, $\omega_n = \sqrt{10/1.4} = \SI{2.67}{rad/s}$ and $\zeta = 7/(2\sqrt{14}) = 0.94$, so $\zeta\omega_n = 2.5$ and $\omega_n\sqrt{1 - \zeta^2} = 0.945$, as found.
\end{example}

\section{Any order, and any input}

\begin{theorem}[Linear equations with constant coefficients]
\label{thm:A-order-n}
Consider $a_n x^{(n)} + \dots + a_1\dot x + a_0 x = 0$ with real coefficients, $a_n \ne 0$, and let its characteristic polynomial $p(s) = a_ns^n + \dots + a_0$ have the distinct roots $s_1, \dots, s_r$ with multiplicities $\mu_1, \dots, \mu_r$.
\begin{enumerate}[label=(\roman*)]
  \item The $n$ functions $t^k e^{s_i t}$, $k = 0, \dots, \mu_i - 1$, are solutions, and every solution is a combination of them with constant coefficients.
  \item For every set of initial values $x(0), \dot x(0), \dots, x^{(n-1)}(0)$ there is exactly one solution.
\end{enumerate}
(Hirsch, Smale and Devaney, 2013, ch.~6.)
\end{theorem}

Three consequences are used constantly in the lessons.
\begin{itemize}
  \item \textbf{Count.} An equation of order $n$ has $n$ poles and needs $n$ initial conditions. The loop with an integral (\lessonref{les:integral}) is of third order; with the motor lag as well, fourth.
  \item \textbf{Superpose.} Because the equation is linear, the sum of two solutions is a solution, and the response to a sum of inputs is the sum of the responses. This is why a step twice as large gives a response twice as large, and why it fails as soon as a motor saturates.
  \item \textbf{Complex roots come in pairs.} The coefficients are real, so if $\sigma + j\omega$ is a root then so is $\sigma - j\omega$, and together they give a real oscillation $e^{\sigma t}\cos(\omega t + \varphi)$.
\end{itemize}

\subsection{Fitting initial conditions}
With distinct roots, $x = \sum_i C_i e^{s_i t}$ and its derivatives at $t = 0$ are $x^{(k)}(0) = \sum_i s_i^k C_i$: $n$ linear equations for the $n$ constants. Examples~\ref{ex:droop-exact} and~\ref{ex:integral-overshoot} solve such a set for three poles.

\subsection{A forced equation}
With an input on the right-hand side,
\begin{equation}\label{eq:A-forced}
  a_n x^{(n)} + \dots + a_0 x = f(t),
\end{equation}
every solution is \emph{one} particular solution plus a solution of the equation without input (the \emph{homogeneous} equation). The reason is one line: the difference of two solutions of \cref{eq:A-forced} solves the homogeneous equation. If all roots have negative real parts, the homogeneous part dies out and the particular solution is what remains — the steady response.

For the inputs that matter in control a particular solution can be guessed:
\begin{center}\small
\renewcommand{\arraystretch}{1.3}
\begin{tabular}{@{}lll@{}}
\toprule
input $f(t)$ & try & result\\
\midrule
constant $f_0$ & a constant & $x = f_0/a_0$ \quad (if $a_0 \ne 0$)\\
$e^{\lambda t}$ & $X e^{\lambda t}$ & $X = 1/p(\lambda)$ \quad (if $p(\lambda) \ne 0$)\\
$\sin\omega t$ & the imaginary part of $X e^{j\omega t}$ & $X = 1/p(j\omega)$\\
\bottomrule
\end{tabular}
\end{center}
The second line contains the other two ($\lambda = 0$ and $\lambda = j\omega$), and it is again the miracle of the exponential: substituting $X e^{\lambda t}$ turns every derivative into a factor $\lambda$, so the left-hand side becomes $p(\lambda)\,X e^{\lambda t}$.\pitfall{If $\lambda$ is itself a root of $p$, the guess fails: $p(\lambda) = 0$. The input then excites a natural motion of the system, the response grows like $t\,e^{\lambda t}$, and this is \emph{resonance}.}

\begin{example}[A sinusoidal force on the P–D loop]
\label{ex:A-forced}
A gust pushes the \SI{1}{kg} drone with the force $\sin\omega t$ newtons at \SI{0.4}{Hz}. The loop is $\ddot x + 7\dot x + 10x = \sin\omega t$. Find the steady motion.

\textbf{Step 1 — the frequency.} $\omega = 2\pi \cdot 0.4 = \SI{2.513}{rad/s}$.

\textbf{Step 2 — the complex problem.} Solve $\ddot z + 7\dot z + 10 z = e^{j\omega t}$ instead, with $z = X e^{j\omega t}$. Its imaginary part will solve the real problem, because the equation is linear with real coefficients.

\textbf{Step 3 — the amplitude.} $p(j\omega) = (j\omega)^2 + 7j\omega + 10 = (10 - 6.32) + 17.59j = 3.68 + 17.59j$, so $X = 1/(3.68 + 17.59j)$.

\textbf{Step 4 — magnitude and angle.} $|p| = \sqrt{3.68^2 + 17.59^2} = 17.97$ and $\angle p = \arctan(17.59/3.68) = 78.2^\circ$. Hence $|X| = 1/17.97 = 0.0556$ and $\angle X = -78.2^\circ$.

\textbf{Step 5 — the answer.} $x(t) = 0.0556\,\sin(\omega t - 78.2^\circ)$: the drone swings by \SI{5.6}{cm} and lags the force by almost a quarter of a period. Example~\ref{ex:challenge-freq} does the same calculation with the integral included.
\end{example}

\section{State space}

An equation of order $n$ can always be rewritten as $n$ equations of first order. Take the third-order loop of \lessonref{les:edge}, $\tau_m m\,\dddot x + m\,\ddot x + K_d\,\dot x + K_p\,x = 0$, and give names to the derivatives: $x_1 = x$, $x_2 = \dot x$, $x_3 = \ddot x$. Then $\dot x_1 = x_2$ and $\dot x_2 = x_3$ by definition, and the equation itself, solved for the highest derivative, gives $\dot x_3$:
\begin{equation}\label{eq:A-companion}
  \frac{\dd}{\dd t}\begin{pmatrix} x_1\\ x_2\\ x_3\end{pmatrix}
  = \begin{pmatrix} 0 & 1 & 0\\ 0 & 0 & 1\\ -\dfrac{K_p}{\tau_m m} & -\dfrac{K_d}{\tau_m m} & -\dfrac{1}{\tau_m}\end{pmatrix}
  \begin{pmatrix} x_1\\ x_2\\ x_3\end{pmatrix} .
\end{equation}
A matrix of this shape — ones above the diagonal, the coefficients in the last row — is a \term{companion matrix}.

\begin{definition}{State-space form}
A linear time-invariant system in \term{state-space form} is
\begin{equation}\label{eq:A-ss}
  \dot{\vx} = A\,\vx + B\,\symbf u, \qquad \symbf y = C\,\vx + D\,\symbf u ,
\end{equation}
with the state $\vx \in \mathbb{R}^n$, the input $\symbf u \in \mathbb{R}^p$, the output $\symbf y \in \mathbb{R}^q$ and constant matrices of fitting sizes. $A$ is the \term{state matrix}, $B$ the input matrix, $C$ the output matrix; $D$ passes the input straight to the output and is zero for every plant in this book.
\end{definition}

\begin{theorem}[The characteristic polynomial]
\label{thm:A-charpoly}
The companion matrix of $s^n + a_{n-1}s^{n-1} + \dots + a_0$ has exactly this polynomial as $\det(s\Id - A)$. For any system $\dot{\vx} = A\vx$ the poles — the exponents $s$ of its solutions $e^{st}\symbf v$ — are the eigenvalues of $A$.
\end{theorem}
\begin{proof}
A solution $e^{st}\symbf v$ with $\symbf v \ne 0$ requires $s\symbf v = A\symbf v$. For the companion matrix this reads $v_2 = s v_1$, $v_3 = s v_2$, \dots, and the last row gives $s\,v_n = -a_0v_1 - \dots - a_{n-1}v_n$. With $v_1 = 1$ the first $n - 1$ equations give $v_k = s^{k-1}$, and the last one becomes the polynomial. So the eigenvalues are its roots, and two polynomials of degree $n$ with leading coefficient one and the same roots (with multiplicity) are equal.
\end{proof}

The proof also gives the eigenvectors of a companion matrix: $(1, s, s^2, \dots)^\T$. They say that in a pure motion $e^{st}$ the velocity is $s$ times the position and the acceleration $s$ times the velocity — which is what differentiating an exponential means.

\begin{example}[The loop with its motors, as a matrix]
\label{ex:A-state}
Write the loop of \lessonref{les:edge} for $K_p = 100$, $K_d = 7$, $m = 1$, $\tau_m = 0.03$ in state form and find its poles.

\textbf{Step 1 — the last row.} $K_p/(\tau_m m) = 3333$, $K_d/(\tau_m m) = 233.3$, $1/\tau_m = 33.3$.

\textbf{Step 2 — the matrix.} $A = \begin{pmatrix} 0 & 1 & 0\\ 0 & 0 & 1\\ -3333 & -233.3 & -33.3\end{pmatrix}$.

\textbf{Step 3 — the polynomial.} By \cref{thm:A-charpoly}, $\det(s\Id - A) = s^3 + 33.3\,s^2 + 233.3\,s + 3333$: the cubic \eqref{eq:edge-cubic} divided by $\tau_m m$.

\textbf{Step 4 — the eigenvalues.} Numerically $-29.25$ and $-2.04 \pm 10.48j$. The pair is the ringing of \cref{fig:edge-runs} at $K_p = 100$: a period of $2\pi/10.48 = \SI{0.60}{s}$ and a decay rate of \SI{2.04}{s^{-1}} (the table of Example~\ref{ex:edge-ringing}, with the D filter and the hold included, has 0.576 and 1.87).
\end{example}

\subsection{The state is not unique}
Any invertible matrix $T$ defines new coordinates $\symbf z = T\vx$, in which
\begin{equation}
  \dot{\symbf z} = T A T^{-1}\,\symbf z + T B\,\symbf u, \qquad \symbf y = C T^{-1}\,\symbf z .
\end{equation}
The same system, described by different numbers. The matrices $A$ and $TAT^{-1}$ are called \emph{similar}; they have the same eigenvalues, because $\det(s\Id - TAT^{-1}) = \det\big(T(s\Id - A)T^{-1}\big) = \det(s\Id - A)$. Poles belong to the system, not to the coordinates. The freedom is used all the time: position and velocity are natural for a drone, but a controller design may prefer coordinates in which $A$ is diagonal, and an estimator may carry extra states of its own.\tip{When two people get different matrices for the same system, compare the eigenvalues of $A$ first. If they agree, the two descriptions probably differ only by a change of coordinates.}

\section{The matrix exponential}

The scalar equation $\dot x = ax$ has the solution $e^{at}x_0$. The vector equation $\dot{\vx} = A\vx$ has the solution $e^{At}\vx_0$, provided the exponential of a matrix is defined by the same power series:
\begin{equation}\label{eq:A-expm}
  e^{At} = \Id + At + \frac{(At)^2}{2!} + \frac{(At)^3}{3!} + \dots
\end{equation}

\begin{theorem}[Properties of the matrix exponential]
\label{thm:A-expm}
For every square matrix $A$ the series \eqref{eq:A-expm} converges for all $t$, and
\begin{enumerate}[label=(\roman*)]
  \item $\frac{\dd}{\dd t}e^{At} = A\,e^{At} = e^{At}A$, and $e^{A\cdot 0} = \Id$;
  \item $e^{A(t + s)} = e^{At}e^{As}$; in particular $e^{At}$ is invertible, with the inverse $e^{-At}$;
  \item if $A\symbf v = \lambda\symbf v$, then $e^{At}\symbf v = e^{\lambda t}\symbf v$;
  \item $e^{TAT^{-1}t} = T\,e^{At}\,T^{-1}$;
  \item $e^{(A + B)t} = e^{At}e^{Bt}$ if $AB = BA$ — and in general \emph{only} then.
\end{enumerate}
\end{theorem}
\begin{proof}
Convergence: the norm of the $k$-th term is at most $(\|A\|\,|t|)^k/k!$, the terms of the series of $e^{\|A\||t|}$. (i)~Differentiate term by term. (ii)~Both sides solve $\dot X = AX$ as functions of $t$ with the same value at $t = 0$, and the solution is unique (\cref{thm:meet-solution}). (iii)~$A^k\symbf v = \lambda^k\symbf v$ in every term. (iv)~$(TAT^{-1})^k = TA^kT^{-1}$ in every term. (v)~With $AB = BA$ the binomial theorem holds for $(A + B)^k$ and the product of the two series can be rearranged as for numbers.
\end{proof}

Three ways to compute $e^{At}$ by hand cover every case in this book.

\paragraph{1. The series stops.} If some power of $A$ is zero, the series is a polynomial. For the double integrator, $A = \begin{psmallmatrix}0 & 1\\ 0 & 0\end{psmallmatrix}$ has $A^2 = 0$, so $e^{At} = \Id + At = \begin{psmallmatrix}1 & t\\ 0 & 1\end{psmallmatrix}$: position grows by velocity times time, velocity stays (\lessonref{les:meet}).

\paragraph{2. Diagonalise.} If $A$ has $n$ independent eigenvectors, collect them as the columns of $V$ and the eigenvalues in $\Lambda = \diag(\lambda_1, \dots, \lambda_n)$. Then $A = V\Lambda V^{-1}$, and by (iv)
\begin{equation}\label{eq:A-diag}
  e^{At} = V\,\diag\!\big(e^{\lambda_1 t}, \dots, e^{\lambda_n t}\big)\,V^{-1} .
\end{equation}

\paragraph{3. One eigenvalue, twice.} A $2 \times 2$ matrix with a double eigenvalue $\lambda$ satisfies $(A - \lambda\Id)^2 = 0$. Write $A = \lambda\Id + N$ with $N = A - \lambda\Id$; the two parts commute, so by (v) and method~1, $e^{At} = e^{\lambda t}(\Id + Nt)$. This is where the factor $t$ of critical damping comes from.

\begin{example}[The exponential of the P–D loop]
\label{ex:A-expm}
Compute $e^{At}$ for the default altitude loop, $A = \begin{psmallmatrix}0 & 1\\ -10 & -7\end{psmallmatrix}$.

\textbf{Step 1 — eigenvalues.} $\det(s\Id - A) = s^2 + 7s + 10 = (s + 2)(s + 5)$: $\lambda_1 = -2$, $\lambda_2 = -5$.

\textbf{Step 2 — eigenvectors.} For a companion matrix they are $(1, \lambda)^\T$: $\symbf v_1 = (1, -2)^\T$, $\symbf v_2 = (1, -5)^\T$.

\textbf{Step 3 — the matrices.} $V = \begin{pmatrix}1 & 1\\ -2 & -5\end{pmatrix}$, $\det V = -3$, $V^{-1} = \dfrac13\begin{pmatrix}5 & 1\\ -2 & -1\end{pmatrix}$.

\textbf{Step 4 — multiply out} \cref{eq:A-diag}:
\begin{equation*}
  e^{At} = \frac13\begin{pmatrix}
    5e^{-2t} - 2e^{-5t} & e^{-2t} - e^{-5t}\\[2pt]
    -10e^{-2t} + 10e^{-5t} & -2e^{-2t} + 5e^{-5t}
  \end{pmatrix} .
\end{equation*}

\textbf{Step 5 — check.} At $t = 0$ it is the identity. Its first column is the motion from $x_0 = 1$, $v_0 = 0$: the position $\tfrac13(5e^{-2t} - 2e^{-5t})$ and, below it, its derivative. The second column is the motion from a unit velocity.
\end{example}

\simfig{tb_phase}{Four matrices, four portraits: the solutions $e^{At}\vx_0$ of $\ddot x + c\,\dot x + k\,x = 0$ in the plane of position and velocity, computed from \cref{eq:A-expm}. With two real poles the curves line up along the eigenvectors (grey). With complex poles they spiral. Without damping they close. With a spring that pushes away, $k < 0$, one eigenvector leads in and the other out: a saddle.}{fig:A-phase}

\subsection{Modes}
Property (iii) has a picture (\cref{fig:A-phase}). If the state starts on an eigenvector, it stays on it and merely shrinks or grows: $\vx(t) = e^{\lambda t}\symbf v$. Such a motion is a \term{mode}. With $n$ independent eigenvectors every initial state is a combination $\vx_0 = \sum c_i\symbf v_i$, and then
\begin{equation}
  \vx(t) = \sum_i c_i\,e^{\lambda_i t}\,\symbf v_i .
\end{equation}
The modes do not interact; each runs at its own rate. In the top-left panel the fast mode ($\lambda = -5$) dies first, and every trajectory ends by sliding in along the slow eigenvector $(1, -2)^\T$: near the end of a response the velocity is always $-2$ times the distance to the setpoint. That is what a \emph{dominant pole} looks like in the plane.

\section{The response to an input}

\Cref{thm:meet-solution} gave the solution of $\dot{\vx} = A\vx + B\symbf u$:
\begin{equation}\label{eq:A-voc}
  \vx(t) = e^{At}\,\vx_0 + \int_0^t e^{A(t - s)}\,B\,\symbf u(s)\,\dd s .
\end{equation}
It is \cref{eq:A-conv1} with matrices, and it is derived the same way: multiply the equation by $e^{-At}$ and recognise $\frac{\dd}{\dd t}\big(e^{-At}\vx\big) = e^{-At}B\symbf u$. The name is \term{variation of constants}. The output is $\symbf y = C\vx$, so for a system at rest at $t = 0$
\begin{equation}\label{eq:A-conv}
  \symbf y(t) = \int_0^t h(t - s)\,\symbf u(s)\,\dd s, \qquad h(t) = C\,e^{At}B .
\end{equation}
The function $h$ is the \term{impulse response}: what the output does after a single, infinitely short push of unit area at $t = 0$. The integral is a \term{convolution}: every past input contributes the impulse response, shifted to its own time and scaled by its size.\recall{The discrete version is \cref{thm:noise-variance}: $z_k = \sum_j h_j\,n_{k-j}$. There $h_j$ was the response of a filter to one unit sample.}

\begin{example}[A step of force, by the formula]
\label{ex:A-step}
The default loop of Example~\ref{ex:A-expm} is at rest at the setpoint when a constant force of \SI{1}{N} appears. Find $x(t)$.

\textbf{Step 1 — the input matrix.} The force enters the acceleration: $B = (0, 1/m)^\T = (0, 1)^\T$, and $u = 1$.

\textbf{Step 2 — the impulse response.} $e^{At}B$ is the second column of $e^{At}$; its first entry, the position, is $h(t) = \tfrac13\,(e^{-2t} - e^{-5t})$.

\textbf{Step 3 — integrate.} With $\vx_0 = 0$ and $u = 1$, $x(t) = \int_0^t h(t - s)\,\dd s = \int_0^t h(\sigma)\,\dd\sigma$:
\begin{equation*}
  x(t) = \frac13\Big[\frac{1 - e^{-2t}}{2} - \frac{1 - e^{-5t}}{5}\Big] = 0.1 - \tfrac16\,e^{-2t} + \tfrac1{15}\,e^{-5t} .
\end{equation*}

\textbf{Step 4 — check.} $x(0) = 0.1 - 0.1667 + 0.0667 = 0$. For $t \to \infty$, $x \to \SI{0.1}{m}$: the droop $F/K_p$ of \lessonref{les:droop}. And the step response is the integral of the impulse response — a general fact, visible in Step~3.
\end{example}

\section{Controllability and observability}

Two questions about \cref{eq:A-ss} decide what feedback can do at all. Can the input move the state wherever we want? And can the output tell us where the state is?

\begin{definition}{Controllable, observable}
The system \eqref{eq:A-ss} is \term{controllable} if for every pair of states $\vx_0$, $\vx_1$ and every time $t_1 > 0$ there is an input that takes the state from $\vx_0$ at $t = 0$ to $\vx_1$ at $t_1$. It is \term{observable} if the initial state $\vx_0$ can be computed from the input and the output over any interval $[0, t_1]$.
\end{definition}

\begin{theorem}[The rank tests]
\label{thm:A-rank}
For a system with $n$ states,
\begin{align*}
  \text{controllable} &\iff \operatorname{rank}\begin{pmatrix} B & AB & \cdots & A^{n-1}B\end{pmatrix} = n ,\\
  \text{observable} &\iff \operatorname{rank}\begin{pmatrix} C^\T & (CA)^\T & \cdots & (CA^{n-1})^\T\end{pmatrix} = n .
\end{align*}
The second matrix is written on its side to save space: its transpose has the rows $C$, $CA$, \dots, $CA^{n-1}$.
(Kálmán, 1960; Kailath, 1980, ch.~2.)
\end{theorem}

Why these matrices? By \cref{eq:A-voc,eq:A-expm} the input reaches the state only through $e^{A(t-s)}B$, a combination of $B$, $AB$, $A^2B$, \dots: $B$ is the direction the input pushes directly, $AB$ where that push is carried by the dynamics, and so on. Powers beyond $A^{n-1}$ add nothing new, because every matrix satisfies its own characteristic polynomial (the Cayley–Hamilton theorem), so $A^n$ is a combination of the lower powers. For observability, differentiate the output with $\symbf u = 0$: $\symbf y = C\vx$, $\dot{\symbf y} = CA\vx$, $\ddot{\symbf y} = CA^2\vx$, \dots. The output and its derivatives at one instant are the rows $C$, $CA$, \dots\ times the state; the state can be solved for exactly when those rows have rank $n$.

\begin{example}[What the drone can be told, and what it can tell]
\label{ex:A-rank}
The L1 drone has $A = \begin{psmallmatrix}0 & 1\\ 0 & 0\end{psmallmatrix}$, $B = \begin{psmallmatrix}0\\ 1/m\end{psmallmatrix}$ and the state (altitude, velocity).

\textbf{Step 1 — controllability.} $AB = (1/m, 0)^\T$, so $(B\ \ AB) = \begin{psmallmatrix}0 & 1/m\\ 1/m & 0\end{psmallmatrix}$: rank 2. One force steers both the altitude and the velocity.

\textbf{Step 2 — an altimeter.} $C = (1\ \ 0)$ and $CA = (0\ \ 1)$: rank 2, observable. The velocity can be recovered from the altitude, by differentiating — which is what the D term does, and what the Kalman filter of Part~II does better.

\textbf{Step 3 — a speed sensor alone.} $C = (0\ \ 1)$ and $CA = (0\ \ 0)$: rank 1, \emph{not} observable. From the velocity one can compute how far the drone has moved, never where it started. Every navigation system that integrates velocities or accelerations has this problem: its position drifts, and only a position measurement can correct it.

\textbf{Step 4 — a system that is not controllable.} Two identical motors driven by the same command: $\dot x_1 = -x_1/\tau + u/\tau$, $\dot x_2 = -x_2/\tau + u/\tau$. Here $B = (1, 1)^\T/\tau$ and $AB = -(1, 1)^\T/\tau^2$: rank 1. The difference $x_1 - x_2$ obeys $\frac{\dd}{\dd t}(x_1 - x_2) = -(x_1 - x_2)/\tau$, with no input in it at all. It cannot be steered. To turn, a multicopter needs motors with separate commands.
\end{example}

The rank tests say yes or no. How \emph{well} a system can be controlled or observed — with how much effort, with how much noise — is a matter of degree, and the tools for it (Gramians, singular values) are in Appendix~I and Lesson~III.9.

\begin{result}[Appendix A at a glance]
\begin{center}\small
\renewcommand{\arraystretch}{1.4}
\begin{tabularx}{\linewidth}{@{}l>{\raggedright\arraybackslash}X@{}}
first order & $\tau\dot x + x = u_0$: \quad $x = u_0 + (x_0 - u_0)\,e^{-t/\tau}$\\
order $n$ & try $e^{st}$; the poles are the roots of the characteristic polynomial; a root of multiplicity $\mu$ contributes $e^{st}, t\,e^{st}, \dots, t^{\mu-1}e^{st}$\\
forcing $e^{\lambda t}$ & steady response $e^{\lambda t}/p(\lambda)$; for $\sin\omega t$ use $\lambda = j\omega$ and take the imaginary part\\
state space & $\dot{\vx} = A\vx + B\symbf u$, $\symbf y = C\vx$; poles $=$ eigenvalues of $A$\\
solution & $\vx(t) = e^{At}\vx_0 + \int_0^t e^{A(t-s)}B\symbf u(s)\,\dd s$\\
$e^{At}$ & $V\diag(e^{\lambda_it})V^{-1}$ if $A = V\Lambda V^{-1}$; \ $\Id + At$ if $A^2 = 0$\\
controllable & $\operatorname{rank}(B\ \ AB\ \cdots\ A^{n-1}B) = n$\\
observable & $\operatorname{rank}(C;\ CA;\ \dots;\ CA^{n-1}) = n$\\
\end{tabularx}
\end{center}
\end{result}

\toolchapter{B}{Polynomials and stability}{Whether the roots lie in the left half-plane can be read from the coefficients. Where they go when a gain changes can be sketched without solving anything.}
\label{app:routh}

\lead{\usedcard{\setuprow{Lesson I.2}{a missing coefficient means instability}\setuprow{Lessons I.5, I.10}{the largest stable $K_i$ and $K_p$}\setuprow{Lesson I.14}{the edge of a fourth-order chain}\setuprow{Lessons I.4, I.5, I.10}{paths of the poles}\setuprow{Part III}{the root locus as a live chart (III.2)}}%
A linear loop is stable exactly when all roots of its characteristic polynomial have negative real parts (\cref{thm:spring-linear}). Computing the roots of a cubic or a quartic by hand is tedious, and with a gain left as a symbol it is hopeless. The two tools of this appendix avoid it. The Routh–Hurwitz criterion decides stability from the coefficients, with symbols allowed. The root locus shows how the roots move as one gain grows.}

\section{Roots and coefficients}

Write a polynomial with its leading coefficient one and factor it over its roots $s_1, \dots, s_n$:
\begin{equation}
  p(s) = s^n + a_{n-1}s^{n-1} + \dots + a_1 s + a_0 = (s - s_1)(s - s_2)\cdots(s - s_n) .
\end{equation}
Multiplying out the right-hand side and comparing coefficients gives \term{Vieta's formulas}. The two that are used most:
\begin{equation}\label{eq:B-vieta}
  s_1 + s_2 + \dots + s_n = -a_{n-1}, \qquad s_1 s_2 \cdots s_n = (-1)^n a_0 .
\end{equation}
The sum of the poles of a loop is fixed by one coefficient. In the loop with the motor lag, $\tau_m m s^3 + m s^2 + K_ds + K_p$, that coefficient is $1/\tau_m$ whatever the gains: \emph{the sum of the three poles is always $-1/\tau_m$}. If a gain pushes one pole to the left, the others must move to the right. This is the algebra behind \enquote{every lag sets a speed limit}.

Two more facts hold for every real polynomial.
\begin{itemize}
  \item Complex roots come in conjugate pairs $\sigma \pm j\omega$, and a pair contributes the real quadratic factor $(s - \sigma)^2 + \omega^2 = s^2 - 2\sigma s + \sigma^2 + \omega^2$.
  \item If all roots have negative real parts, the polynomial is a product of factors $(s + a)$ and $(s^2 + bs + c)$ with positive $a$, $b$, $c$, so \emph{all its coefficients are positive}. A missing or negative coefficient proves instability at a glance (Step~4 of Example~\ref{ex:spring-cubic}).
\end{itemize}
The converse of the second fact holds up to degree two and fails from degree three on: $s^3 + s^2 + s + 6 = (s + 2)(s^2 - s + 3)$ has positive coefficients and two roots at $+0.5 \pm 1.66j$.

\subsection{A slow root, quickly}
When one root is much closer to the origin than the others, it can be estimated from the last two coefficients. Near $s = 0$ the high powers are negligible and $p(s) \approx a_1 s + a_0$, so
\begin{equation}\label{eq:B-slow}
  s_{slow} \approx -\frac{a_0}{a_1}, \qquad\text{and, one step better,}\qquad s_{slow} \approx -\frac{a_0}{a_1}\Big(1 + \frac{a_0\,a_2}{a_1^2}\Big) .
\end{equation}
The second form keeps the $s^2$ term and inserts the first estimate into it. For the loop of \lessonref{les:droop}, $s^3 + 7s^2 + 10s + 0.8$: $-a_0/a_1 = -0.080$, which is the quasi-static $-K_i/K_p$, and the refined value is $-0.080\,(1 + 0.8 \cdot 7/100) = -0.0845$. The exact root is $-0.0850$.

\section{The Routh–Hurwitz criterion}\index{Routh–Hurwitz criterion|(}

\begin{definition}{Hurwitz polynomial}
A real polynomial is called \term{Hurwitz} (or \emph{stable}) if all its roots have negative real parts.
\end{definition}

\begin{theorem}[Routh–Hurwitz, degrees two to four]
\label{thm:B-low}
Let the leading coefficient be positive. The polynomial is Hurwitz if and only if all its coefficients are positive and, in addition,
\begin{center}
\renewcommand{\arraystretch}{1.35}
\begin{tabular}{@{}ll@{}}
degree 2, \ $a_2s^2 + a_1s + a_0$: & nothing more;\\
degree 3, \ $a_3s^3 + a_2s^2 + a_1s + a_0$: & $a_2a_1 > a_3a_0$;\\
degree 4, \ $a_4s^4 + a_3s^3 + a_2s^2 + a_1s + a_0$: & $a_1\,(a_3a_2 - a_4a_1) > a_3^2\,a_0$.
\end{tabular}
\end{center}
\end{theorem}
\begin{proof}
Degrees two and three are \cref{thm:spring-routh}. For degree four, divide by $a_4$ and let the polynomial be $s^4 + a_3s^3 + a_2s^2 + a_1s + a_0$. Its roots move continuously with the coefficients, so a root can pass from the left half-plane to the right only through the imaginary axis. At $s = j\omega$ the real part is $\omega^4 - a_2\omega^2 + a_0$ and the imaginary part $\omega\,(a_1 - a_3\omega^2)$. Both vanish for some $\omega \ne 0$ exactly when $\omega^2 = a_1/a_3$ and $a_1^2 - a_1a_2a_3 + a_0a_3^2 = 0$, that is, when $a_1(a_3a_2 - a_1) = a_3^2a_0$; a root at $s = 0$ needs $a_0 = 0$. So within the set of positive coefficients the number of right-half-plane roots can change only across the surface $a_1(a_3a_2 - a_1) = a_3^2a_0$. On the side where the inequality holds lies $(s + 1)^4 = s^4 + 4s^3 + 6s^2 + 4s + 1$, with $4 \cdot (24 - 4) = 80 > 16$ and all roots at $-1$; on the other side lies $s^4 + s^3 + s^2 + s + 1$, whose roots are four of the fifth roots of unity, two of them in the right half-plane. Each side is connected, so the count is constant on it.
\end{proof}

The proof shows something useful in itself: \emph{at the edge of stability the roots on the imaginary axis are at $\omega^2 = a_1/a_3$} (degrees three and four). The lessons use this to find the frequency at which a loop oscillates at its limit (Examples~\ref{ex:integral-limit}, \ref{ex:edge-cubic} and~\ref{ex:inv-quartic}).

\subsection{Any degree: the Routh array}
For higher degrees the conditions are organised in a table. Write the coefficients in two rows, alternating, and build each further row from the two above it:
\begin{equation}\label{eq:B-array}
\begin{array}{c|cccc}
  s^n & a_n & a_{n-2} & a_{n-4} & \cdots\\
  s^{n-1} & a_{n-1} & a_{n-3} & a_{n-5} & \cdots\\
  s^{n-2} & b_1 & b_2 & b_3 & \cdots\\
  s^{n-3} & c_1 & c_2 & \cdots & \\
  \vdots & \vdots & & &
\end{array}
\qquad\quad
\begin{aligned}
  b_1 &= \frac{a_{n-1}a_{n-2} - a_n a_{n-3}}{a_{n-1}},\\[2pt]
  b_2 &= \frac{a_{n-1}a_{n-4} - a_n a_{n-5}}{a_{n-1}},\\[2pt]
  c_1 &= \frac{b_1 a_{n-3} - a_{n-1} b_2}{b_1}, \quad\dots
\end{aligned}
\end{equation}
Each entry is a $2 \times 2$ cross product of the first column of the two rows above with the column to its right, divided by the first entry of the row directly above. Missing entries count as zero. The table has $n + 1$ rows.

\begin{theorem}[Routh]
\label{thm:B-routh}
If no entry of the first column of the array \eqref{eq:B-array} is zero, the number of roots with positive real parts equals the number of changes of sign down the first column. In particular the polynomial is Hurwitz if and only if all entries of the first column have the same sign. (Routh, 1877; Gantmacher, 1959, ch.~XV.)
\end{theorem}

\begin{example}[The cascade of Lesson I.14, by the array]
\label{ex:B-array}
The chain of three loops with the velocity integral has the characteristic polynomial $s^4 + K s^3 + 2.2K s^2 + 3.24K s + 0.72K$ (\cref{eq:inv-quartic}, $K = K_{att}$). Test it for $K = 3$ and $K = 1.5$.

\textbf{Step 1 — $K = 3$: the first two rows.} The coefficients are $1,\ 3,\ 6.6,\ 9.72,\ 2.16$:
\begin{equation*}
\begin{array}{c|ccc}
  s^4 & 1 & 6.6 & 2.16\\
  s^3 & 3 & 9.72 & \\
\end{array}
\end{equation*}

\textbf{Step 2 — the $s^2$ row.} $b_1 = (3 \cdot 6.6 - 1 \cdot 9.72)/3 = 3.36$ and $b_2 = (3 \cdot 2.16 - 1 \cdot 0)/3 = 2.16$.

\textbf{Step 3 — the $s^1$ row.} $c_1 = (3.36 \cdot 9.72 - 3 \cdot 2.16)/3.36 = 7.79$.

\textbf{Step 4 — the $s^0$ row.} $d_1 = (7.79 \cdot 2.16 - 3.36 \cdot 0)/7.79 = 2.16$. The last entry of the first column is always $a_0$.

\textbf{Step 5 — read the column.} $1,\ 3,\ 3.36,\ 7.79,\ 2.16$: no change of sign. Stable. (The roots are $-0.42 \pm 2.03j$, $-1.90$, $-0.26$.)

\textbf{Step 6 — $K = 1.5$.} Coefficients $1,\ 1.5,\ 3.3,\ 4.86,\ 1.08$. Then $b_1 = (1.5 \cdot 3.3 - 4.86)/1.5 = 0.06$, $b_2 = 1.08$, and $c_1 = (0.06 \cdot 4.86 - 1.5 \cdot 1.08)/0.06 = -22.1$. The column is $1,\ 1.5,\ 0.06,\ -22.1,\ 1.08$: two changes of sign, from $+$ to $-$ and back. Two roots in the right half-plane — the pair $+0.043 \pm 1.76j$ that makes the drone of \cref{fig:inversion-predict} swing for ever.
\end{example}

\subsection{When a whole row vanishes}
At the edge of stability the array breaks down in a telling way: a row becomes zero. The row above it then holds the coefficients of a polynomial in $s^2$ — the \term{auxiliary polynomial} — whose roots are the roots of $p$ that lie symmetrically about the origin, and in particular the pair on the imaginary axis.

In Example~\ref{ex:B-array} with $K$ as a symbol the $s^2$ row is $b_1 = 2.2K - 3.24$, $b_2 = 0.72K$, and the $s^1$ entry vanishes when $b_1 \cdot 3.24K = K \cdot 0.72K$, i.e.\ $K = 1.638$. The auxiliary polynomial is $b_1 s^2 + b_2 = 0.364\,s^2 + 1.179$, with the roots $\pm j\sqrt{1.179/0.364} = \pm 1.80j$: the edge and its frequency, as in Example~\ref{ex:inv-quartic}.\tip{With a gain as a symbol, do not compute the whole array. Write the entries of the first column as functions of the gain and ask where each changes sign. The smallest such gain is the edge.}

\subsection{The same test as determinants}
Hurwitz found the criterion independently in 1895 in another form, which is the one to use in a proof or in a program. For $p(s) = a_ns^n + \dots + a_0$ with $a_n > 0$ build the $n \times n$ \term{Hurwitz matrix}
\begin{equation}
  H = \begin{pmatrix}
    a_{n-1} & a_{n-3} & a_{n-5} & \cdots\\
    a_n & a_{n-2} & a_{n-4} & \cdots\\
    0 & a_{n-1} & a_{n-3} & \cdots\\
    0 & a_n & a_{n-2} & \cdots\\
    \vdots & & & \ddots
  \end{pmatrix}.
\end{equation}
The polynomial is Hurwitz if and only if the determinants of the upper-left $1 \times 1$, $2 \times 2$, \dots, $n \times n$ blocks of $H$ are all positive. For $n = 3$ they are $a_2$, $a_2a_1 - a_3a_0$ and $a_0$ times the second: the conditions of \cref{thm:B-low}. The entries of Routh's first column are the ratios of successive determinants, which is why the two tests agree.\index{Routh–Hurwitz criterion|)}

\section{Stability with a margin}

\enquote{All roots in the left half-plane} is the least one can ask. Often one wants more: all roots to the left of the line $\operatorname{Re}s = -\sigma$, so that every motion decays at least as fast as $e^{-\sigma t}$. The Routh–Hurwitz test answers that too, after a shift of the variable.

\begin{result}[Testing a decay rate]
All roots of $p(s)$ satisfy $\operatorname{Re}s < -\sigma$ if and only if the polynomial $q(w) = p(w - \sigma)$ is Hurwitz.
\end{result}

The reason: $s = w - \sigma$ moves every root to the right by $\sigma$, so $\operatorname{Re}w < 0$ means $\operatorname{Re}s < -\sigma$.

\begin{example}[Does the default loop decay faster than $e^{-t}$?]
\label{ex:B-shift}
The P–D loop has $p(s) = s^2 + 7s + 10$. Test $\sigma = 1$ and $\sigma = 2.5$.

\textbf{Step 1 — shift by 1.} $q(w) = (w - 1)^2 + 7(w - 1) + 10 = w^2 + 5w + 4$. All coefficients positive: Hurwitz. Every motion of the loop decays faster than $e^{-t}$.

\textbf{Step 2 — shift by 2.5.} $q(w) = (w - 2.5)^2 + 7(w - 2.5) + 10 = w^2 + 2w - 1.25$. A negative coefficient: not Hurwitz. Some motion decays more slowly than $e^{-2.5t}$.

\textbf{Step 3 — compare.} The roots are $-2$ and $-5$. The slowest, $-2$, lies between the two lines, as the test says.
\end{example}

\section{The root locus}\index{root locus|(}

The Routh test says \emph{whether} a loop is stable for a given gain. The root locus shows \emph{how} the poles move as the gain grows, and it can be sketched by rules.

Let the characteristic equation depend on one gain $k \ge 0$ in the form
\begin{equation}\label{eq:B-locus}
  D(s) + k\,N(s) = 0 ,
\end{equation}
with polynomials $D$ of degree $n_p$ and $N$ of degree $n_z \le n_p$, both with leading coefficient one. The roots of $D$ are called the (open-loop) \emph{poles} $p_i$, the roots of $N$ the \emph{zeros} $z_i$. Every characteristic equation of this book can be brought into this form for any one of its gains: collect the terms with the gain in $kN$ and the rest in $D$.

\begin{definition}{Root locus}
The \term{root locus} of \cref{eq:B-locus} is the set of all points $s$ of the complex plane that are roots for some $k \ge 0$. It consists of $n_p$ curves, the \emph{branches}, each traced by one root as $k$ runs from $0$ to $\infty$.
\end{definition}

Rewrite \cref{eq:B-locus} as $N(s)/D(s) = -1/k$. A complex number equals a negative real number when its angle is $180^\circ$ (up to full turns) and its magnitude fits. With the factored forms this gives two conditions:
\begin{align}
  \text{angle:} \quad & \sum_i \angle(s - z_i) - \sum_i \angle(s - p_i) = 180^\circ + \ell \cdot 360^\circ, \label{eq:B-angle}\\
  \text{magnitude:} \quad & k = \frac{\prod_i |s - p_i|}{\prod_i |s - z_i|} . \label{eq:B-mag}
\end{align}
The angle condition alone decides whether a point is on the locus; the magnitude condition then tells for which gain. All the rules below are consequences of \cref{eq:B-angle}.

\begin{theorem}[Rules for sketching a root locus]
\label{thm:B-locus}
For \cref{eq:B-locus} with $k$ from $0$ to $\infty$:
\begin{enumerate}[label=(\roman*)]
  \item \emph{Start and end.} The branches start at the poles ($k = 0$). As $k \to \infty$, $n_z$ of them end at the zeros and $n_p - n_z$ go to infinity.
  \item \emph{Symmetry.} The locus is symmetric about the real axis.
  \item \emph{Real axis.} A point of the real axis is on the locus exactly when the number of real poles and zeros to its right is odd.
  \item \emph{Asymptotes.} The branches that go to infinity approach straight lines through the point
  \begin{equation*}
    \sigma_a = \frac{\sum p_i - \sum z_i}{n_p - n_z}
    \qquad\text{at the angles}\qquad \frac{(2\ell + 1)\,180^\circ}{n_p - n_z},
  \end{equation*}
  with $\ell = 0, 1, \dots, n_p - n_z - 1$.
  \item \emph{Meeting points.} Branches meet and part (on the real axis: break away or break in) at roots of $N\,D' - N'D = 0$ that lie on the locus.
  \item \emph{Crossing the imaginary axis.} The gain and the frequency of a crossing follow from the Routh–Hurwitz edge condition.
\end{enumerate}
(Evans, 1948; Franklin, Powell and Emami-Naeini, 2019, ch.~5.)
\end{theorem}
\begin{proof}[Why]
(i)~For $k = 0$ the equation is $D = 0$; dividing by $k$ and letting $k \to \infty$ it becomes $N = 0$. (ii)~Real coefficients. (iii)~For a real point $s$, a real pole or zero to its right contributes the angle $180^\circ$, one to its left $0^\circ$, and a complex pair contributes angles that cancel; \cref{eq:B-angle} then needs an odd count. (iv)~Far away all poles and zeros are seen under nearly the same angle $\theta$, so \cref{eq:B-angle} reads $(n_z - n_p)\,\theta = 180^\circ + \ell\,360^\circ$; the centre follows from the two highest coefficients of $D + kN$ by \cref{eq:B-vieta}. (v)~Where branches meet, \cref{eq:B-locus} has a double root, so its derivative $D' + kN' = 0$ vanishes as well; eliminate $k$. (vi)~\Cref{thm:B-low}.
\end{proof}

\begin{example}[Raising $K_p$ with a lagging motor]
\label{ex:B-locus-kp}
Sketch the path of the poles of \lessonref{les:edge} as $K_p$ grows, with $K_d = 7$, $m = 1$, $\tau_m = 0.03$.

\textbf{Step 1 — the form.} $\tau_m m s^3 + m s^2 + K_d s + K_p = 0$. Divide by $\tau_m m$: $D(s) = s^3 + 33.3\,s^2 + 233.3\,s$, $N(s) = 1$ and $k = K_p/(\tau_m m) = 33.3\,K_p$.

\textbf{Step 2 — poles and zeros.} $D = s\,(s^2 + 33.3s + 233.3) = s\,(s + 10)(s + 23.3)$: poles at $0$, $-10$, $-23.3$. No zeros: $n_p - n_z = 3$.

\textbf{Step 3 — the real axis.} On the locus: the segment between $0$ and $-10$ (one pole to the right) and everything left of $-23.3$ (three).

\textbf{Step 4 — asymptotes.} Centre $\sigma_a = (0 - 10 - 23.3)/3 = -11.1$; angles $60^\circ$, $180^\circ$, $300^\circ$. One branch runs left along the real axis. Two go up and down to the right, and must cross the imaginary axis.

\textbf{Step 5 — the breakaway.} With $N = 1$ rule (v) is $D'(s) = 3s^2 + 66.7s + 233.3 = 0$: $s = -4.35$ or $-17.9$. Only $-4.35$ is on the locus. The gain there, from \cref{eq:B-mag}: $k = 4.35 \cdot 5.65 \cdot 18.98 = 466$, i.e.\ $K_p = 14.0$.

\textbf{Step 6 — the crossing.} \Cref{thm:B-low}: $a_2a_1 = a_3a_0$ at $33.3 \cdot 233.3 = k$, so $k = 7778$, $K_p = 233$, at $\omega = \sqrt{a_1/a_3} = \sqrt{233.3} = \SI{15.3}{rad/s}$.

\textbf{Step 7 — read the sketch} (\cref{fig:B-locus}, left). Up to $K_p = 14$ the two slow poles are real and approach each other: the loop is overdamped. At 14 they meet — critical damping, a little above the $K_d^2/4m = 12.25$ of the ideal loop. Beyond it they leave the axis and bend to the right, slowly at first, and cross into the right half-plane at $K_p = 233$. Meanwhile the third pole, the motor, moves left: the sum of the three stays at $-33.3$.
\end{example}

\simfig{tb_locus}{Two root loci of the altitude loop with a motor lag of \SI{30}{ms}, computed from the characteristic polynomial. Left: $K_p$ grows while $K_d = 7$ stays (Example~\ref{ex:B-locus-kp}). Right: $K_p$ and $K_d$ grow together in the ratio $10 : 7$ (Example~\ref{ex:B-locus-ratio}). Crosses: open-loop poles; circle: the zero; dots: the default tune; dashed: the asymptotes.}{fig:B-locus}

\begin{example}[Raising $K_p$ and $K_d$ together]
\label{ex:B-locus-ratio}
Now keep the ratio $K_p/K_d = 10/7$ and scale both gains. How does the picture change?

\textbf{Step 1 — the form.} $\tau_m m s^3 + m s^2 + K_d\,(s + K_p/K_d) = 0$: $D = s^3 + 33.3\,s^2 = s^2(s + 33.3)$, $N = s + 1.43$, $k = K_d/(\tau_m m) = 33.3\,K_d$.

\textbf{Step 2 — poles and zeros.} A double pole at $0$ (the drone) and one at $-33.3$ (the motor). One zero at $-1.43$. $n_p - n_z = 2$.

\textbf{Step 3 — the real axis.} Between $-33.3$ and $-1.43$ (three poles and zeros to the right). Not between $-1.43$ and $0$, where the count is two.

\textbf{Step 4 — asymptotes.} Centre $\sigma_a = (0 + 0 - 33.3 + 1.43)/2 = -15.9$, angles $\pm 90^\circ$: two vertical lines, \emph{in the left half-plane}.

\textbf{Step 5 — meeting points.} $N D' - N' D = (s + 1.43)(3s^2 + 66.7s) - (s^3 + 33.3s^2) = 0$ has the roots $0$, $-3.01$ and $-15.8$. The two poles from the origin circle round the zero and land on the real axis at $-3.01$, for $K_d = 5.2$; a pair leaves it again at $-15.8$, for $K_d = 9.1$.

\textbf{Step 6 — read the sketch} (\cref{fig:B-locus}, right). For $K_d$ between 5.2 and 9.1 all three poles are real; the default $K_d = 7$ lies in this window, with poles at $-1.93$, $-7.11$ and $-24.3$. For larger gains a pair climbs the asymptote at $-15.9$: it oscillates faster and faster and keeps its decay rate. No gain makes this loop unstable. The edge of \lessonref{les:edge} exists because $K_p$ was raised \emph{alone}, which moves the zero $-K_p/K_d$ to the left and with it the asymptotes to the right.
\end{example}

\begin{keyidea}[Zeros attract, and the excess decides]
Poles move towards zeros as the gain grows. The poles that find no zero leave along the asymptotes, and their number $n_p - n_z$ decides the fate of the loop at high gain: one or two can stay in the left half-plane for ever; three or more cannot (\cref{thm:edge-asymptotes}). A controller improves a loop by adding zeros in good places — which is what the D term does.
\end{keyidea}\index{root locus|)}

\section{The same question for sampled loops}

A sampled loop is stable when the roots of its characteristic polynomial in $z$ lie \emph{inside the unit circle} (\cref{thm:slow-sampled}). The Routh–Hurwitz test does not apply directly, but a change of variable makes it apply: the map
\begin{equation}\label{eq:B-bilinear}
  z = \frac{1 + w}{1 - w}, \qquad w = \frac{z - 1}{z + 1}
\end{equation}
sends the inside of the unit circle in the $z$-plane to the left half of the $w$-plane. (For $w = j v$ on the imaginary axis, $|1 + jv| = |1 - jv|$, so $|z| = 1$; and $w = -1$ gives $z = 0$.)

\begin{theorem}[Stability of a second-order recursion]
\label{thm:B-jury}
Both roots of $z^2 + a_1 z + a_0$ lie inside the unit circle if and only if
\begin{equation*}
  |a_0| < 1, \qquad 1 + a_1 + a_0 > 0, \qquad 1 - a_1 + a_0 > 0 .
\end{equation*}
\end{theorem}
\begin{proof}
Insert \cref{eq:B-bilinear} and multiply by $(1 - w)^2$: $(1 + w)^2 + a_1(1 + w)(1 - w) + a_0(1 - w)^2 = (1 - a_1 + a_0)\,w^2 + 2\,(1 - a_0)\,w + (1 + a_1 + a_0)$. By \cref{thm:B-low} this is Hurwitz exactly when its three coefficients have the same sign. They cannot all be negative, because the first and the last add up to $2 + 2a_0$ while the middle one is $2 - 2a_0$, and the two cannot both be negative. So all three are positive; the first and last give the two inequalities, and with them the middle one gives $a_0 < 1$ while their sum gives $a_0 > -1$.
\end{proof}

The three conditions are three lines in the plane of $(a_1, a_0)$, and they enclose a triangle: the stability region of every second-order digital filter and loop. Appendix~D uses it to find the slowest sample rate of a P–D controller.

\begin{result}[Appendix B at a glance]
\begin{center}\small
\renewcommand{\arraystretch}{1.4}
\begin{tabularx}{\linewidth}{@{}l>{\raggedright\arraybackslash}X@{}}
necessary & all coefficients present and of one sign\\
degree 2 & that is also sufficient\\
degree 3 & $a_2a_1 > a_3a_0$; at the edge $\omega^2 = a_1/a_3$\\
degree 4 & $a_1(a_3a_2 - a_4a_1) > a_3^2a_0$; at the edge $\omega^2 = a_1/a_3$\\
any degree & Routh array: sign changes in the first column $=$ roots in the right half-plane\\
decay rate $\sigma$ & test $p(w - \sigma)$\\
slow root & $s \approx -a_0/a_1$\\
sum of roots & $-a_{n-1}/a_n$\\
root locus & $D + kN = 0$: from the poles to the zeros and to $n_p - n_z$ asymptotes at $(2\ell + 1)180^\circ/(n_p - n_z)$ through $(\sum p_i - \sum z_i)/(n_p - n_z)$\\
sampled loop & $z^2 + a_1z + a_0$ stable iff $|a_0| < 1$ and $1 \pm a_1 + a_0 > 0$\\
\end{tabularx}
\end{center}
\end{result}

